\section{AIME I 2011}
        \begin{enumerate}
        	\item (\textit{AIME I 2011}) Jar $A$ contains four liters of a solution that is $45\%$ acid. Jar $B$ contains five liters of a solution that is $48\%$ acid. Jar $C$ contains one liter of a solution that is $k\%$ acid. From jar $C$, $\frac{m}{n}$ liters of the solution is added to jar $A$, and the remainder of the solution in jar $C$ is added to jar B. At the end both jar $A$ and jar $B$ contain solutions that are $50\%$ acid. Given that $m$ and $n$ are relatively prime positive integers, find $k + m + n$.


 
 \textbf{Answer: }085



\textbf{Solution 1}

Jar A contains $\frac{11}{5}$ liters of water, and $\frac{9}{5}$ liters of acid; jar B contains $\frac{13}{5}$ liters of water and $\frac{12}{5}$ liters of acid.

 The gap between the amount of water and acid in the first jar, $\frac{2}{5}$, is double that of the gap in the second jar, $\frac{1}{5}$. Therefore, we must add twice as much of jar C into the jar $A$ over jar $B$. So, we must add $\frac{2}{3}$ of jar C into jar $A$, so $m = 2, n=3$.

 Since jar C contains $1$ liter of solution, we are adding $\frac{2}{3}$ of a liter of solution to jar $A$. In order to close the gap between water and acid, there must be $\frac{2}{5}$ more liters of acid than liters of water in these $\frac{2}{3}$ liters of solution. So, in the $\frac{2}{3}$ liters of solution, there are $\frac{2}{15}$ liters of water, and $\frac{8}{15}$ liters of acid.  So, 80\% of the $\frac{2}{3}$ sample is acid, so overall, in jar C, 80\% of the sample is acid.

 Therefore, our answer is $80 + 2 + 3 = \boxed{85}$.

 





\textbf{Solution 2}

There are $\frac{45}{100}(4)=\frac{9}{5}$ L of acid in Jar A. There are $\frac{48}{100}(5)=\frac{12}{5}$ L of acid in Jar B. And there are $\frac{k}{100}$ L of acid in Jar C. After transferring the solutions from jar C, there will be
 $4+\frac{m}{n}$ L of solution in Jar A and $\frac{9}{5}+\frac{k}{100}\cdot\frac{m}{n}$ L of acid in Jar A.

 $6-\frac{m}{n}$ L of solution in Jar B and $\frac{12}{5}+\frac{k}{100}\cdot \left(1-\frac{m}{n}\right)=\frac{12}{5}+\frac{k}{100}-\frac{mk}{100n}$ of acid in Jar B.
Since the solutions are 50\% acid, we can multiply the amount of acid for each jar by 2, then equate them to the amount of solution.

 \[\frac{18}{5}+\frac{km}{50n}=4+\frac{m}{n}\]
 \[\frac{24}{5}-\frac{km}{50n}+\frac{k}{50}=6-\frac{m}{n}\]Add the equations to get
 \[\frac{42}{5}+\frac{k}{50}=10\] Solving gives $k=80$.
If we substitute back in the original equation we get $\frac{m}{n}=\frac{2}{3}$ so $3m=2n$. Since $m$ and $n$ are relatively prime, $m=2$ and $n=3$. Thus $k+m+n=80+2+3=\boxed{085}$.

 





\textbf{Solution 3}

One might cleverly change the content of both Jars. 

 Since the end result of both Jars are $50\%$ acid, we can turn Jar A into a 1 gallon liquid with $50\%-4(5\%) = 30\%$ acid 

 and Jar B into 1 gallon liquid with $50\%-5(2\%) =40\%$ acid.

 Now, since Jar A and Jar B contain the same amount of liquid, twice as much liquid will be pour into Jar A than Jar B, so $\dfrac{2}{3}$ of Jar C will be pour into Jar A.

 Thus, $m=2$ and $n=3$.

 $\dfrac{30\% + \frac{2}{3} \cdot k\%}{\frac{5}{3}} = 50\%$

 Solving for $k$ yields $k=80$

 So the answer is $80+2+3 = \boxed{085}$

 





\textbf{Solution 4}

One may first combine all three jars in to a single container.  That container will have $10$ liters of liquid, and it should be $50\%$ acidic.  Thus there must be $5$ liters of acid.  

 Jar A contained $45\% \cdot 4L$, or $1.8L$ of acid, and jar B $48\% \cdot 5L$ or $2.4L$.  Solving for the amount of acid in jar C, $k = (5 - 2.4 - 1.8) = .8$, or $80\%$.

 Once one knows that the jar C is $80\%$ acid, use solution 1 to figure out m and n for $k+m+n=80+2+3=\boxed{085}$.

 




	\item (\textit{AIME I 2011}) In rectangle $ABCD$, $AB = 12$ and $BC = 10$.  Points $E$ and $F$ lie inside rectangle $ABCD$ so that $BE = 9$, $DF = 8$, $\overline{BE} \parallel \overline{DF}$, $\overline{EF} \parallel \overline{AB}$, and line $BE$ intersects segment $\overline{AD}$.  The length $EF$ can be expressed in the form $m \sqrt{n} - p$, where $m$, $n$, and $p$ are positive integers and $n$ is not divisible by the square of any prime.  Find $m + n + p$.


 
 \textbf{Answer: }036



\textbf{Solution 1}

Let us call the point where $\overline{EF}$ intersects $\overline{AD}$ point $G$, and the point where $\overline{EF}$ intersects $\overline{BC}$ point $H$. Since angles $FHB$ and $EGA$ are both right angles, and angles $BEF$ and $DFE$ are congruent due to parallelism, right triangles $BHE$ and $DGF$ are similar. This implies that $\frac{BH}{GD} = \frac{9}{8}$. Since $BC=10$, $BH+GD=BH+HC=BC=10$. ($HC$ is the same as $GD$ because they are opposite sides of a rectangle.) Now, we have a system:

 $\frac{BH}{GD}=\frac{9}8$

 $BH+GD=10$

 Solving this system (easiest by substitution), we get that:

 $BH=\frac{90}{17}$

 $GD=\frac{80}{17}$

 Using the Pythagorean Theorem, we can solve for the remaining sides of the two right triangles:

 $\sqrt{9^2-\left(\frac{90}{17}\right)^2}$ and $\sqrt{8^2-\left(\frac{80}{17}\right)^2}$

 Notice that adding these two sides would give us twelve plus the overlap $EF$. This means that:

 $EF= \sqrt{9^2-\left(\frac{90}{17}\right)^2}+\sqrt{8^2-\left(\frac{80}{17}\right)^2}-12=3\sqrt{21}-12$

 Since $21$ isn't divisible by any perfect square, our answer is:

 $3+21+12=\boxed{36}$

 





\textbf{Solution 2}

Extend lines $BE$ and $CD$ to meet at point $G$.Draw the altitude $GH$ from point $G$ to line $BA$ extended.

 $GE=DF=8,$$GB=17$

 In right $\bigtriangleup GHB$, $GH=10$, $GB=17$, thus by Pythagoras Theorem we have:$HB=\sqrt{17^2-10^2}=3\sqrt{21}$

 $HA=EF=3\sqrt{21}-12$

 Thus our answer is:$3+21+12=\boxed{36}$

 


 





\textbf{Solution 3 (Similar triangles, Algebra)}

We notice that since $\overline{BE}||\overline{DF}$, $\overline{EF}$ is the diagonal of a rectangle. Now, let us extend the width lines to intersect with $\overline{BE}$ and $\overline{DF}$, respectively, to form this rectangle. Let us call the length of $\overline{EF}$ $d$, the perpendicular distance between $\overline{EF}$ and $\overline{AD}$ $k$, and the perpendicular distance between $\overline{EF}$ and $\overline{CD}$ $x$. We now can begin the similar triangles. When drawing the diagram (where $E$ is closer to $\overline{AD}$ than $F$), we put the similar right triangles in the same position so that we can begin solving for our variables and finding ratios. Since all of these triangles are similar, we find that $\frac{dx}{8}=\frac{10d-dx}{9}$, which solves for $x=\frac{80}{17}$. Completing the same thing for $k$, we see $\frac{12d-dk}{9}=\frac{dk+d^2}{8}$, which solves for $k=\frac{96-9d}{17}$. We are now ready to find both the length and the width of the rectangle. We find these to be $\frac{10d}{17}$ and $\frac{12d+d^2}{17}$. Now let us use the Pythagorean to solve for $d$. We square the length and the width and multiply both sides by 289 (the denominator of the LHS) to reach the equation $(12+d^2)^2+100d^2=289d^2$. Expanding and simplifying, we find that $d^4+24d^3=45d^2$. Dividing both sides by $d^2$ and moving everything to the LHS, we see that $d^2+24d-45=0$. Applying the quadratic formula, $d=\frac{-24\pm \sqrt{756}}{2}=-12\pm \sqrt{189}$. Since $d>0$, $d=\sqrt{189}-12=3\sqrt{21}-12$. Finally, $m+n+p=3+21+12=\boxed{036}$. -Gideontz

 




	\item (\textit{AIME I 2011}) Let $L$ be the line with slope $\frac{5}{12}$ that contains the point $A = (24,-1)$, and let $M$ be the line perpendicular to line $L$ that contains the point $B = (5,6)$.  The original coordinate axes are erased, and line $L$ is made the $x$-axis and line $M$ the $y$-axis.  In the new coordinate system, point $A$ is on the positive $x$-axis, and point $B$ is on the positive $y$-axis.  The point $P$ with coordinates $(-14,27)$ in the original system has coordinates $(\alpha,\beta)$ in the new coordinate system.  Find $\alpha + \beta$.


 
 \textbf{Answer: }031



\textbf{Solution}

Given that $L$ has slope $\frac{5}{12}$ and contains the point $A=(24,-1)$, we may write the point-slope equation for $L$ as $y+1=\frac{5}{12}(x-24)$.Since $M$ is perpendicular to $L$ and  contains the point $B=(5,6)$, we have that the slope of $M$ is $-\frac{12}{5}$, and consequently that the point-slope equation for $M$ is $y-6=-\frac{12}{5}(x-5)$.

 


 Converting both equations to the form $0=Ax+By+C$, we have that $L$ has the equation $0=5x-12y-132$ and that $M$ has the equation $0=12x+5y-90$.Applying the point-to-line distance formula, $\frac{|Ax+By+C|}{\sqrt{A^2+B^2}}$, to point $P$ and lines $L$ and $M$, we find that the distance from $P$ to $L$ and $M$ are $\frac{526}{13}$ and $\frac{123}{13}$, respectively. 

 


 Since $A$ and $B$ lie on the positive axes of the shifted coordinate plane, we may show by graphing the given system that point P will lie in the second quadrant in the new coordinate system. Therefore, the $x$-coordinate of $P$ is negative, and is therefore $-\frac{123}{13}$; similarly, the $y$-coordinate of $P$ is positive and is therefore $\frac{526}{13}$.

 Thus, we have that $\alpha=-\frac{123}{13}$ and that $\beta=\frac{526}{13}$. It follows that $\alpha+\beta=-\frac{123}{13}+\frac{526}{13}=\frac{403}{13}=\boxed{031}$.

 





\textbf{Solution 2 (alternate bash)}

The equations for the axes are $\frac{5}{12} (x-24) = y+1$ and $-\frac{12}{5}(x-5) = y - 6$. We can solve the system to find that they intersect at the point $\left( \frac{1740}{169},\frac{-1134}{169} \right)$

 The unit basis vectors of our new axes are $\begin{pmatrix} 12/13 \\ 5/13 \end{pmatrix}$ and $\begin{pmatrix} -5/13 \\ 12/13 \end{pmatrix}$ for the $x$ and $y$ axes respectively (taking into account which direction is positive). 

 Then, we solve the following system for $\alpha$ and $\beta$ :

 
 \[\alpha \begin{pmatrix} 12/13 \\ 5/13 \end{pmatrix} + \beta \begin{pmatrix} -5/13 \\ 12/13 \end{pmatrix} + \begin{pmatrix} 1740/169 \\ -1134/169 \end{pmatrix} = \begin{pmatrix} -14 \\ 27 \end{pmatrix}\]

 Painful bashing gives $\alpha = -\frac{123}{13}$ and $\beta = \frac{526}{13}$. Adding gives $\alpha + \beta = \frac{403}{13} = \boxed{031}$

 We can also attempt to manipulate the system of equations to solve for $\alpha + \beta$, and avoid solving for $\alpha$ and $\beta$ separately.

 \textbf{Possibly a solution 3} First, find where the new positive x and y axis are by comparing the points A and B to the line it is not on (higher or lower). Next, just find the new quadrant the point is in from the x and y axis. Finally, to find the distance from the point to the axis, we can just use both slopes to find both projections of the line onto the axis. The intersection of the projection line and the axis's previous line would give us information on the new point's coordinate points. 

 





\textbf{Solution 4 (Complex Palne)}

As in solution 2, find where the 2 lines intersect:$\left( \frac{1740}{169},\frac{-1134}{169} \right)$then use complex numbers. Let $p=-14+27i$ denote the complex number representing $P$. Let $p'=\alpha+\beta i$ denote the desired complex number. Then:
 \[p'=\left(-14-\frac{1740}{169}+27i+\frac{1134}{169}\right)\div\text{arg}(12+5i)\]
 \[p'=\left(-14-\frac{1740}{169}+27i+\frac{1134}{169}\right)\div\frac{12+5i}{13}\]
 \[p'=\left(-14-\frac{1740}{169}+27i+\frac{1134}{169}\right)\cdot\frac{12-5i}{13}\]Note that to find $\alpha+\beta$, one does not need to bash all the way, rather, after mutplying the fraction with the bracket, one can add all the coefficients to find the answer of $\boxed{031}$.

 




	\item (\textit{AIME I 2011}) In triangle $ABC$, $AB = 125$, $AC = 117$, and $BC = 120$.  The angle bisector of angle $A$ intersects $\overline{BC}$ at point $L$, and the angle bisector of angle $B$ intersects $\overline{AC}$ at point $K$.  Let $M$ and $N$ be the feet of the perpendiculars from $C$ to $\overline{BK}$ and $\overline{AL}$, respectively.  Find $MN$.


 
 \textbf{Answer: }056



\textbf{Solution 1}

Extend ${CM}$ and ${CN}$ such that they intersect line ${AB}$ at points $P$ and $Q$, respectively. 
                    \begin{figure}[H]
                        \centering
                        \includegraphics[width=0.5\textwidth]{Solution/AIME_I_2011_4_Solution_Figure_1.png}
                    \end{figure}
                    Since ${BM}$ is the angle bisector of angle $B$ and ${CM}$ is perpendicular to ${BM}$, $\triangle BCP$ must be an isoceles triangle, so $BP=BC=120$, and $M$ is the midpoint of ${CP}$. For the same reason, $AQ=AC=117$, and $N$ is the midpoint of ${CQ}$.Hence $MN=\frac 12 PQ$. Since 
 \[PQ=BP+AQ-AB=120+117-125=112,\] so $MN=\boxed{056}$.

 





\textbf{Solution 2}

Let $I$ be the incenter of $ABC$. Since $I$ lies on $BM$ and $AN$, $IM \perp MC$ and $IN \perp NC$, so $\angle IMC + \angle INC = 180^\circ$. This means that $CMIN$ is a cyclic quadrilateral. By the Law of Sines, $\frac{MN}{\sin \angle MIN} = \frac{2R}{\sin \angle CMI} = 2R = CI$, where $R$ is the radius of the circumcircle of $CMIN$. Since $\sin \angle MIN = \sin \angle BIA = \sin (90^\circ + \frac 12 \angle BCA) = \cos \frac 12 \angle BCA = \cos \angle BCI$, we have that $MN = CI \cdot \sin \angle MIN = CI \cdot \cos \angle BCI$. Letting $H$ be the point of contact of the incircle of $ABC$ with side $BC$, we have $MN = CI \cdot \cos \angle BCI = CI \cdot \frac{CH}{CI} = CH$. Thus, $MN = s - AB = \frac{117+120-125}{2}=\boxed{056}$.

 





\textbf{Solution 3 (Bash)}

Project $I$ onto $AC$ and $BC$ as $D$ and $E$. $ID$ and $IE$ are both in-radii of $\triangle ABC$ so we get right triangles with legs $r$ (the in-radius length) and $s - c = 56$. Since $IC$ is the hypotenuse for the 4 triangles ($\triangle INC, \triangle IMC, \triangle IDC,$ and $\triangle IEC$), $C, D, M, I, N, E$ are con-cyclic on a circle we shall denote as $\omega$ which is also the circumcircle of $\triangle CMN$ and $\triangle CDE$. To find $MN$, we can use the Law of Cosines on $\angle MON \implies MN^2 = 2R^2(1 - \cos{2\angle MCN})$ where $O$ is the center of $\omega$. Now, the circumradius $R$ can be found with Pythagorean Theorem with $\triangle CDI$ or $\triangle CEI$: $r^2 + 56^2 = (2R)^2$. To find $r$, we can use the formula $rs = [ABC]$ and by Heron's, $[ABC] = \sqrt{181 \cdot 61 \cdot 56 \cdot 64} \implies r = \sqrt{\frac{61 \cdot 56 \cdot 64}{181}} \implies 2R^2 = \frac{393120}{181}$. To find $\angle MCN$, we can find $\angle MIN$ since $\angle MCN = 180 - \angle MIN$. $\angle MIN = \angle MIC + \angle NIC = 180 - \angle BIC + 180 - \angle AIC = 180 - (180 - \frac{\angle A + \angle C}{2}) + 180 - (180 - \frac{\angle B + \angle C}{2}) = \frac{\angle A + \angle B + \angle C}{2} + \frac{\angle C}{2}$. Thus, $\angle MCN = 180 - \frac{\angle A + \angle B + \angle C}{2} - \frac{\angle C}{2}$ and since $\angle A + \angle B + \angle C = 180$, we have $\angle A + \angle B + \angle C - \frac{\angle A + \angle B + \angle C}{2} - \frac{\angle C}{2} = \frac{\angle A + \angle B}{2}$. Plugging this into our Law of Cosines (LoC) formula gives $MN^2 = 2R^2(1 - \cos{\angle A + \angle B}) = 2R^2(1 + \cos{\angle C})$. To find $\cos{\angle C}$, we use LoC on $\triangle ABC \implies \cos{\angle C} = \frac{120^2 + 117^2 - 125^2}{2 \cdot 117 \cdot 120} = \frac{41 \cdot 19}{117 \cdot 15}$. Our formula now becomes $MN^2 = \frac{393120}{181} + \frac{2534}{15 \cdot 117}$. After simplifying, we get $MN^2 = 3136 \implies MN = \boxed{056}$.

 --lucasxia01

 





\textbf{Solution 4}

Because $\angle CMI = \angle CNI = 90$, $CMIN$ is cyclic. 

 Applying Ptolemy's theorem on CMIN:

 $CN \cdot MI+CM \cdot IN=CI \cdot MN$

 $CI^2(\cos \angle ICN \sin \angle ICM + \cos \angle ICM \sin \angle ICN) = CI \cdot MN$

 $MN = CI \sin \angle MCN$ by sine angle addition formula.

 $\angle MCN = 180 - \angle MIN = 90 - \angle BCI$. 

 Let $H$ be where the incircle touches $BC$, then $CI \cos \angle BCI = CH = \frac{a+b-c}{2}$.$a=120, b=117, c=125$, for a final answer of $\boxed{056}$.

 Note: This is similar to Solution 2 after the first four lines

 





\textbf{Solution 5 (Trig Bash)}

Applying \href{https://artofproblemsolving.com/wiki/index.php?title=Ptolemy\%27s\_Theorem}{Ptolemy's Theorem} on the cyclic quadrilateral $MINC$, we find that

 $MI\cdot CN + IN\cdot MC = MN\cdot IC$.

 $\angle CIN=\frac{\alpha+\gamma}{2}$ and $\angle MIC=\frac{\beta+\gamma}{2}$ by the Exterior Angle Theorem, so from properties of sine and cosine, we can find that

 $MI=IC\cdot\cos\left(\frac{\beta+\gamma}{2}\right),$ $MC=IC\cdot\sin\left(\frac{\beta+\gamma}{2}\right),$ $IN=IC\cdot\cos\left(\frac{\alpha+\gamma}{2}\right),$ $CN=IC\cdot\sin\left(\frac{\alpha+\gamma}{2}\right).$

 Plugging in the values and simplifying results in $MN = IC\cdot\sin\left(\frac{\alpha+\beta+2\gamma}{2}\right)$ by the angle-addition identity $\sin(A+B)=\sin(A)\cos(B)+\sin(B)\cos(A)$.

 Before we continue, we would like to simplify the value in the sine function. We see that $\frac{\alpha+\beta+2\gamma}{2}=\frac{\gamma}{2}+\frac{\alpha+\beta+\gamma}{2}=\frac{\gamma}{2}+90$. Using the fact that $\cos(A)=\sin(90-A)$ results in

 $\sin\left(\frac{\alpha+\beta+2\gamma}{2}\right)=\sin\left(90+\frac{\gamma}{2}\right)=\sin\left(90-(-\frac{\gamma}{2})\right)=\cos\left(-\frac{\gamma}{2}\right)=\cos\left(\frac{\gamma}{2}\right).$

 How do we simplify $IC$? Well, we can perform the \href{https://artofproblemsolving.com/wiki/index.php?title=Law\_of\_Sines}{Law of Sines} on triangle $AIC$. This results in:

 $\frac{AC}{\sin(\angle AIC)}=\frac{IC}{\sin\left(\frac{\alpha}{2}\right)}$

 The value of $\angle AIC$ is $\frac{\alpha+2\beta+\gamma}{2}$ by the Exterior Angle Theorem on $\triangle ABI$, so the value of $\sin(\angle AIC)$ is equivalent to the value of $\cos\left(\frac{\beta}{2}\right)$ by a similar argument as above. Then rearranging yields $IC = b\cdot\frac{\sin\left(\frac{\alpha}{2}\right)}{\cos\left(\frac{\beta}{2}\right)}$.

 Going back to the previous formula $MN = IC\cdot\sin\left(\frac{\alpha+\beta+2\gamma}{2}\right)$ and substituting values yields:

 $MN = b\cdot\frac{\sin\left(\frac{\alpha}{2}\right)\cos\left(\frac{\gamma}{2}\right)}{\cos\left(\frac{\beta}{2}\right)}$.

 Finally, using the formulae $\sin\left(\frac{\alpha}{2}\right)=\sqrt{\frac{(s-b)(s-c)}{bc}}$ and $\cos\left(\frac{\alpha}{2}\right)=\sqrt{\frac{s(s-a)}{bc}}$ (where $s$ is half the perimeter of the triangle), we reach our final value:

 $MN = b\cdot\frac{\sqrt{\frac{(s-b)(s-c)}{bc}}\cdot\sqrt{\frac{s(s-c)}{ab}}}{\sqrt{\frac{s(s-b)}{ac}}}$

 $=\frac{\sqrt{s(s-b)(s-c)^2}}{\sqrt{s(s-b)}}$

 $=\sqrt{(s-c)^2}$

 $=s-c$

 $=181-125$

 $=\boxed{056}.$

 




	\item (\textit{AIME I 2011}) The vertices of a regular nonagon (9-sided polygon) are to be labeled with the digits 1 through 9 in such a way that the sum of the numbers on every three consecutive vertices is a multiple of 3.  Two acceptable arrangements are considered to be indistinguishable if one can be obtained from the other by rotating the nonagon in the plane.  Find the number of distinguishable acceptable arrangements.


 
 \textbf{Answer: }144



\textbf{Solution 1}

First, we determine which possible combinations of digits $1$ through $9$ will yield sums that are multiples of $3$. It is simplest to do this by looking at each of the digits $\bmod{3}$.

 We see that the numbers $1, 4,$ and $7$ are congruent to $1 \pmod{3}$, that the numbers $2, 5,$ and $8$ are congruent to $2 \pmod{3}$, and that the numbers $3, 6,$ and $9$ are congruent to $0 \pmod{3}$. In order for a sum of three of these numbers to be a multiple of three, the mod $3$ sum must be congruent to $0$. Quick inspection reveals that the only possible combinations are $0+0+0, 1+1+1, 2+2+2,$ and $0+1+2$. However, every set of three consecutive vertices must sum to a multiple of three, so using any of $0+0+0, 1+1+1$, or $2+2+2$ would cause an adjacent sum to include exactly $2$ digits with the same $\bmod{3}$ value, and this is an unacceptable arrangement. Thus the only possible groupings are composed of three digits congruent to three different $\bmod{3}$ values. 

 We see also that there are two possible arrangements for these trios on the nonagon: a digit congruent to $1 \pmod{3}$ can be located counterclockwise of a digit congruent to $0$ and clockwise of a digit congruent to $2 \pmod{3}$, or the reverse can be true. 

 We set the first digit as $3$ avoid overcounting rotations, so we have one option as a choice for the first digit.  The other two $0 \pmod{3}$ numbers can be arranged in $2!=2$ ways.  The three $1 \pmod{3}$ and three $2 \pmod{3}$ can both be arranged in $3!=6$ ways.  Therefore, the desired result is $2(2 \times 6 \times 6)=\boxed{144}$.

 





\textbf{Solution 2}

Notice that there are three triplets of congruent integers $\mod 3$: $(1,4,7),$ $(2,5,8),$ and $(3,6,9).$ There are $3!$ ways to order each of the triplets individually and $3!$ ways to order the triplets as a group (see solution 1). Rotations are indistinguishable, so there are $(3!)^4/9=\boxed{144}$ total arrangements.

 





\textbf{Solution 3 (Cyclic Residues and Equivalencies)}

Essentially, if we sum up each three consecutive vertices, we should see equivalencies modulo 3. To approach this, label the vertices $a, b, c, \dots, i$ in that order. Then, cyclical residues give

 
 \[a+b+c \equiv b+c+d \equiv c+d+e \equiv \dots \equiv h+i+a \equiv 0 \pmod{3}\].

 From the cyclic orientation of all these residues, it follows that $a \equiv d \equiv g \pmod{3}$, $b \equiv e \equiv h \pmod{3}$, and $c \equiv f \equiv i \pmod{3}$. Residue three classes for each give three ordered pairs congruent modulo 3, $(1,4,7),$ $(2,5,8),$ and $(3,6,9)$. Each of these pairs can be given to each of the ordered $(a,d,g)$, $(b,e,h)$, $(c,f,i)$. Each pair has $3!$ orderings, meaning there are $(3!)^3 = 216$ total combinations with one placement. Let $A,B,C =$ $(a,d,g)$, $(b,e,h)$, $(c,f,i)$, and an addition multiplicity of 6 is attributed to the previous 216 combinations, giving 1296 total. Each of these is not accounting for the multiplicity by the 9 rotations of the nonagon, and therefore there is a grand total of $\frac{1296}{9} = \boxed{144}$ distinct arrangements.

 




	\item (\textit{AIME I 2011}) Suppose that a parabola has vertex $\left(\frac{1}{4},-\frac{9}{8}\right)$ and equation $y = ax^2 + bx + c$, where $a > 0$ and $a + b + c$ is an integer.  The minimum possible value of $a$ can be written in the form $\frac{p}{q}$, where $p$ and $q$ are relatively prime positive integers.  Find $p + q$.


 
 \textbf{Answer: }011



\textbf{Solution 1}

If the vertex is at $\left(\frac{1}{4}, -\frac{9}{8}\right)$, the equation of the parabola can be expressed in the form 
 \[y=a\left(x-\frac{1}{4}\right)^2-\frac{9}{8}.\]Expanding, we find that 
 \[y=a\left(x^2-\frac{x}{2}+\frac{1}{16}\right)-\frac{9}{8},\] and 
 \[y=ax^2-\frac{ax}{2}+\frac{a}{16}-\frac{9}{8}.\] From the problem, we know that the parabola can be expressed in the form $y=ax^2+bx+c$, where $a+b+c$ is an integer. From the above equation, we can conclude that $a=a$, $-\frac{a}{2}=b$, and $\frac{a}{16}-\frac{9}{8}=c$. Adding up all of these gives us 
 \[\frac{9a-18}{16}=a+b+c.\] We know that $a+b+c$ is an integer, so $9a-18$ must be divisible by $16$.  Let $9a=z$. If ${z-18}\equiv {0} \pmod{16}$, then ${z}\equiv {2} \pmod{16}$. Therefore, if $9a=2$, $a=\frac{2}{9}$. Adding up gives us $2+9=\boxed{011}$

 





\textbf{Solution 2}

Complete the square. Since $a>0$, the parabola must be facing upwards. $a+b+c=\text{integer}$ means that $f(1)$ must be an integer. The function can be recasted into $a\left(x-\frac{1}{4}\right)^2-\frac{9}{8}$ because the vertex determines the axis of symmetry and the critical value of the parabola. The least integer greater than $-\frac{9}{8}$ is $-1$. So the $y$-coordinate must change by $\frac{1}{8}$ and the $x$-coordinate must change by $1-\frac{1}{4}=\frac{3}{4}$. Thus, $a\left(\frac{3}{4}\right)^2=\frac{1}{8}\implies \frac{9a}{16}=\frac{1}{8}\implies a=\frac{2}{9}$. So $2+9=\boxed{011}$.

 





\textbf{Solution 3}

To do this, we can use the formula for the minimum (or maximum) value of the $x$ coordinate at a vertex of a parabola, $-\frac{b}{2a}$ and equate this to $\frac{1}{4}$. Solving, we get $-\frac{a}{2}=b$. Enter $x=\frac{1}{4}$ to get $-\frac{9}{8}=\frac{a}{16}+\frac{b}{4}+c=-\frac{a}{16}+c$ so $c=\frac{a-18}{16}$. This means that $\frac{9a-18}{16}\in \mathbb{Z}$ so the minimum of $a>0$ is when the fraction equals -1, so $a=\frac{2}{9}$. Therefore, $p+q=2+9=\boxed{011}$.-Gideontz

 





\textbf{Solution 4}

Write this as $a\left( x- \frac 14 \right)^2 - \frac 98$. Since $a+b+c$ is equal to the value of this expression when you plug $x=1$ in, we just need $\frac{9a}{16}- \frac 98$ to be an integer. Since $a>0$, we also have $\frac{9a}{16}>0$ which means $\frac{9a}{16}- \frac 98 > -\frac{9}{8}$. The least possible value of $a$ is when this is equal to $-1$, or $a=\frac 29$, which gives answer $11$.

 -bobthegod78, krwang, Simplest14

 





\textbf{Solution 5 (You don't remember conic section formulae)}

Take the derivative to get that the vertex is at $2ax+b=0$ and note that this implies $\frac{1}{2} \cdot a = -b$ and proceed with any of the solutions above.

 
$\textbf{Appendix to Solution 5 (You don't remember differential calculus)}$

 Note that the quadratic formula for finding roots of parabolas is $x=\frac{-b\pm\sqrt{b^2-4ac}}{2a}$, so if you average the two x-values of the roots the $\pm\sqrt{b^2-4ac}$ part will cancel out and leave you with $x=-\frac{b}{2a}$. Then proceed as above.

 


 





\textbf{Solution 6 (No Words)}

$y = ax^2 + bx + c$

 $y=a\left(x^2+\frac{b}{a}x\right)+c$

 $y=a\left(x^2+\frac{b}{a}x+\left(\frac{b}{2a}\right)^2\right)-a\left(\frac{b}{2a}\right)^2+c$

 $y=a\left(x+\frac{b}{2a}\right)^2-\frac{b^2}{4a}+c$.

 


 $\frac{b}{2a}=-\frac{1}{4} \implies b=-\frac{1}{2}a$.

 
$- \frac{b^2}{4a} + c= - \frac{9}{8}$

 $c=\frac{b^2}{4a}- \frac{9}{8} =\frac{b}{2} \cdot \frac{b}{2a} - \frac{9}{8} = -\frac{b}{8}-\frac{9}{8} = \frac{a}{16}-\frac{9}{8}$.

 


 $a+b+c=a-\frac{1}{2}a+\frac{a}{16}-\frac{9}{8} = \frac{16a-8a+a}{16}-\frac{9}{8} = \frac{9a-18}{16} \in \mathbb{Z}$.

 $16 \mid 9(a-2) \implies a-2=b \cdot \frac{16}{9}, b \in \mathbb{Z}$.

 $a=\frac{16b+18}{9} \implies a = \ldots, -\frac{14}{9}, \frac{2}{9}, \frac{18}{9}, \frac{34}{9}, \ldots$.

 $a>0 \implies a_{min} = \frac{2}{9} \implies 2+9 = \boxed{11}$.

 




	\item (\textit{AIME I 2011}) Find the number of positive integers $m$ for which there exist nonnegative integers $x_0$, $x_1$ , $\dots$ , $x_{2011}$ such that

 \[m^{x_0} = \sum_{k = 1}^{2011} m^{x_k}.\]


 
 \textbf{Answer: }016



\textbf{Solution 1}

$m^{x_0}= m^{x_1} +m^{x_2} + .... + m^{x_{2011}}$. Now, divide by $m^{x_0}$ to get $1= m^{x_1-x_0} +m^{x_2-x_0} + .... + m^{x_{2011}-x_0}$. Notice that since we can choose all nonnegative $x_0,...,x_{2011}$, we can make $x_n-x_0$ whatever we desire. WLOG, let $x_0\geq...\geq x_{2011}$ and let $a_n=x_n-x_0$.  Notice that, also, $m^{a_{2011}}$ doesn't matter if we are able to make $m^{a_1} +m^{a_2} + .... + m^{a_{2010}}$ equal to $1-\left(\frac{1}{m}\right)^x$ for any power of $x$. Consider $m=2$. We can achieve a sum of $1-\left(\frac{1}{2}\right)^x$ by doing $\frac{1}{2}+\frac{1}{4}+...$ (the "simplest" sequence). If we don't have $\frac{1}{2}$, to compensate, we need $2\cdot \frac{1}{4}$'s. Now, let's try to generalize. The "simplest" sequence is having $\frac{1}{m}$ $m-1$ times, $\frac{1}{m^2}$ $m-1$ times, $\ldots$. To make other sequences, we can split $m-1$ $\frac{1}{m^i}$s into $m(m-1)$ $\text{ }\frac{1}{m^{i+1}}$s since $(m-1)\cdot\frac{1}{m^{i}} = m(m-1)\cdot\frac{1}{m^{i+1}}$. Since we want $2010$ terms, we have $\sum$ $(m-1)\cdot m^x=2010$. However, since we can set $x$ to be anything we want (including 0), all we care about is that $(m-1)\text{ }|\text{ } 2010$ which happens $\boxed{016}$ times.

 





\textbf{Solution 2}

Let 
 \[P(m) = m^{x_0} - m^{x_1} -m^{x_2} - .... - m^{x_{2011}}\]. The problem then becomes finding the number of positive integer roots $m$ for which $P(m) = 0$ and $x_0, x_1, ..., x_{2011}$ are nonnegative integers. We plug in $m = 1$ and see that 
 \[P(1) = 1 - 1 - 1... -1 = 1-2011 = -2010\] Now, we can say that 
 \[P(m) = (m-1)Q(m) - 2010\] for some polynomial $Q(m)$ with integer coefficients. Then if $P(m) = 0$, $(m-1)Q(m) = 2010$. Thus, if $P(m) = 0$, then $m-1 | 2010$ .Now, we need to show that 
 \[m^{x_{0}}=\sum_{k = 1}^{2011}m^{x_{k}}\forall m-1|2010\] We try with the first few $m$ that satisfy this.For $m = 2$, we see we can satisfy this if $x_0 = 2010$, $x_1 = 2009$, $x_2 = 2008$, $\cdots$ , $x_{2008} = 2$, $x_{2009} = 1$, $x_{2010} = 0$, $x_{2011} = 0$, because 
 \[2^{2009} + 2^{2008} + \cdots + 2^1 + 2^0 +2^ 0 = 2^{2009} + 2^{2008} + \cdots + 2^1 + 2^1 = \cdots\] (based on the idea $2^n + 2^n = 2^{n+1}$, leading to a chain of substitutions of this kind) 
 \[= 2^{2009} + 2^{2008} + 2^{2008} = 2^{2009} + 2^{2009} = 2^{2010}\]. Thus $2$ is a possible value of $m$. For other values, for example $m = 3$, we can use the same strategy, with 
 \[x_{2011} = x_{2010} = x_{2009} = 0\], 
 \[x_{2008} = x_{2007} = 1\], 
 \[x_{2006} = x_{2005} = 2\], 
 \[\cdots\], 
 \[x_2 = x_1 = 1004\] and 
 \[x_0 = 1005\], because 
 \[3^0 + 3^0 + 3^0 +3^1+3^1+3^2+3^2+\cdots+3^{1004} +3^{1004} = 3^1+3^1+3^1+3^2+3^2+\cdots+3^{1004} +3^{1004} = 3^2+3^2+3^2+\cdots+3^{1004} +3^{1004}\]
 \[=\cdots = 3^{1004} +3^{1004}+3^{1004} = 3^{1005}\]It's clearly seen we can use the same strategy for all $m-1 |2010$. We count all positive $m$ satisfying $m-1 |2010$, and see there are $\boxed{016}$

 





\textbf{Solution 3}

One notices that $m-1 \mid 2010$  if and only if there exist non-negative integers $x_0,x_1,\ldots,x_{2011}$ such that $m^{x_0} = \sum_{k=1}^{2011}m^{x_k}$.

 To prove the forward case, we proceed by directly finding $x_0,x_1,\ldots,x_{2011}$. Suppose $m$ is an integer such that $m^{x_0} = \sum_{k=1}^{2011}m^{x_k}$. We will count how many $x_k = 0$, how many $x_k = 1$, etc. Suppose the number of $x_k = 0$ is non-zero. Then, there must be at least $m$ such $x_k$ since $m$ divides all the remaining terms, so $m$ must also divide the sum of all the $m^0$ terms. Thus, if we let $x_k = 0$ for $k = 1,2,\ldots,m$, we have,
 \[m^{x_0} = m + \sum_{k=m+1}^{2011}m^{x_k}.\]Well clearly, $m^{x_0}$ is greater than $m$, so $m^2 \mid m^{x_0}$. $m^2$ will also divide every term, $m^{x_k}$, where $x_k \geq 2$. So, all the terms, $m^{x_k}$, where $x_k < 2$ must sum to a multiple of $m^2$. If there are exactly $m$ terms where $x_k = 0$, then we must have at least $m-1$ terms where $x_k = 1$. Suppose there are exactly $m-1$ such terms and $x_k = 1$ for $k = m+1,m+2,2m-1$. Now, we have,
 \[m^{x_0} = m^2 + \sum_{k=2m}^{2011}m^{x_k}.\]One can repeat this process for successive powers of $m$ until the number of terms reaches 2011. Since there are $m + j(m-1)$ terms after the $j$th power, we will only hit exactly 2011 terms if $m-1$ is a factor of 2010. To see this, 
 \[m+j(m-1) = 2011 \Rightarrow m-1+j(m-1) = 2010 \Rightarrow (m-1)(j+1) = 2010.\]Thus, when $j = 2010/(m-1) - 1$ (which is an integer since $m-1 \mid 2010$ by assumption, there are exactly 2011 terms. To see that these terms sum to a power of $m$, we realize that the sum is a geometric series:
 \[1 + (m-1) + (m-1)m+(m-1)m^2 + \cdots + (m-1)m^j = 1+(m-1)\frac{m^{j+1}-1}{m-1} = m^{j+1}.\]Thus, we have found a solution for the case $m-1 \mid 2010$.

 Now, for the reverse case, we use the formula 
 \[x^k-1 = (x-1)(x^{k-1}+x^{k-2}+\cdots+1).\] Suppose $m^{x_0} = \sum_{k=1}^{2011}m^{x^k}$ has a solution. Subtract 2011 from both sides to get 
 \[m^{x_0}-1-2010 = \sum_{k=1}^{2011}(m^{x^k}-1).\] Now apply the formula to get 
 \[(m-1)a_0-2010 = \sum_{k=1}^{2011}[(m-1)a_k],\] where $a_k$ are some integers. Rearranging this equation, we find 
 \[(m-1)A = 2010,\] where $A = a_0 - \sum_{k=1}^{2011}a_k$. Thus, if $m$ is a solution, then $m-1 \mid 2010$. 

 So, there is one positive integer solution corresponding to each factor of 2010. Since $2010 = 2\cdot 3\cdot 5\cdot 67$, the number of solutions is $2^4 = \boxed{016}$.

 





\textbf{Solution 4 (for noobs like me)}

The problem is basically asking how many integers $m$ have a power that can be expressed as the sum of 2011 other powers of $m$ (not necessarily distinct). Notice that $2+2+4+8+16=32$, $3+3+3+9+9+27+27+81+81=243$, and $4+4+4+4+16+16+16+64+64+64+256+256+256=1024$. Thus, we can safely assume that the equation $2011 = (m-1)x + m$ must have an integer solution $x$. To find the number of $m$-values that allow the aforementioned equation to have an integer solution, we can subtract 1 from the constant $m$ to make the equation equal a friendlier number, $2010$, instead of the ugly prime number $2011$: $2010 = (m-1)x+(m-1)$. Factor the equation and we get $2010 = (m-1)(x+1)$. The number of values of $m-1$ that allow $x+1$ to be an integer is quite obviously the number of factors of $2010$. Factoring $2010$, we obtain $2010 = 2 \times 3 \times 5 \times 67$, so the number of positive integers $m$ that satisfy the required condition is $2^4 = \boxed{016}$.

 -fidgetboss\_4000

 





\textbf{Solution 5}

First of all, note that the nonnegative integer condition really does not matter, since even if we have a nonnegative power, there is always a power of $m$ we can multiply to get to non-negative powers. Now we see that our problem is just a matter of m-chopping blocks. What is meant by $m$-chopping is taking an existing block of say $m^{k}$ and turning it into $m$ blocks of $m^{k-1}$. This process increases the total number of blocks by $m-1$ per chop.The problem wants us to find the number of positive integers $m$ where some number of chops will turn $1$ block into $2011$ such blocks, thus increasing the total amount by $2010= 2 \cdot 3 \cdot 5 \cdot 67$. Thus $m-1 | 2010$, and a cursory check on extreme cases will confirm that there are indeed $\boxed{016}$ possible $m$s.

 





\textbf{Solution 6 (fast, easy)}

The problem is basically saying that we want to find $2011$ powers of $m$ that add together to equal another power of $m$. If we had a power of $m$, $m^n$, then to get to $m^{n+1}$ or $m \cdot (m^n)$, we have to add $m^n$, $m-1$times. Then when we are at $m^{n+1}$, to get to $m^{n+2}$, it is similar. So we have to have $m^{\text{some number}} = m^n + (m-1){m^n} + (m-1)(m^{n+1})\dots$. This expression has $1 + (m-1) + (m-1) + \dots = 1 + k(m-1)$ terms, and the number of terms must be equal to $2011$ (as stated in the problem). Then $1+k(m-1) = 2011$, and we get the equation $k(m-1) = 2010 = 2\cdot3\cdot5\cdot67$. $2010$ has $16$ factors, and setting $m-1$ equal to each of these factors makes $\boxed{016}$ values of $m$.

 





\textbf{Solution 7 (Quick Modular Arithmetic)}

Take$\pmod {m-1}$ of both sides
 \[m^{x_0} \equiv 1 \pmod {m-1}\]
 \[\sum_{k = 1}^{2011} m^{x_k} \equiv 2011 \pmod {m-1}\]Setting both of these to be congruent to eachother, we get:
 \[1 \equiv 2011 \pmod {m-1}\]Subtract $1$ from both sides to get:
 \[0 \equiv 2010 \pmod {m-1}\]$m-1$ must be a factor of $2010$.

 $2010$ has $16$ factors, so there are $\boxed{16}$ values of $m$.

 





\textbf{Solution 8 (Polynomial Construction)}

Let $f(a)$ be a polynomial in terms of $a$, listed below, which has roots at the values $m$ that satisfy the problem condition: 
 \[f(a) = a^{x_0} - \sum_{k=1}^{2011} a^{x_k}\] It is evident to see that the polynomial has integer coefficients, which will be used later in the solution.

 $f(1) = 1 - 2011 = -2010$. We use the fact that polynomials $P(x)$ with integer coefficients satisfy the condition that $(a-b)|f(a)-f(b)$ for some integers $a$ and $b$ to conclude that 
 \[(m-1)|f(m)-f(1)\] Since $m$ is a solution of the polynomial, we have that $f(m) = 0$. Therefore, we have that 
 \[(m-1)|2010\] Therefore, we arrive at the condition that $m-1$ must be a factor of $2010$.

 Since $2010=2\cdot 3\cdot 5\cdot 67$, $2010$ has $16$ factors, so there are $\boxed{016}$ possible values of $m$. 

 -Vivdax

 \textbf{An Olympiad Problem that's (almost) the exact same (and it came before 2011, MAA)} \href{https://artofproblemsolving.com/community/c6h84155p486903}{https://artofproblemsolving.com/community/c6h84155p486903}2001 Austrian-Polish Math Individual Competition \#1~MSC

 




	\item (\textit{AIME I 2011}) In triangle $ABC$, $BC = 23$, $CA = 27$, and $AB = 30$.  Points $V$ and $W$ are on $\overline{AC}$ with $V$ on $\overline{AW}$, points $X$ and $Y$ are on $\overline{BC}$ with $X$ on $\overline{CY}$, and points $Z$ and $U$ are on $\overline{AB}$ with $Z$ on $\overline{BU}$.  In addition, the points are positioned so that $\overline{UV} \parallel \overline{BC}$, $\overline{WX} \parallel \overline{AB}$, and $\overline{YZ} \parallel \overline{CA}$.  Right angle folds are then made along $\overline{UV}$, $\overline{WX}$, and $\overline{YZ}$.  The resulting figure is placed on a level floor to make a table with triangular legs.  Let $h$ be the maximum possible height of a table constructed from triangle $ABC$ whose top is parallel to the floor.  Then $h$ can be written in the form $\frac{k \sqrt{m}}{n}$, where $k$ and $n$ are relatively prime positive integers and $m$ is a positive integer that is not divisible by the square of any prime.  Find $k + m + n$.


 
                    \begin{figure}[H]
                        \centering
                        \includegraphics[width=0.5\textwidth]{Problem/AIME_I_2011_8_Problem_Figure_1.png}
                    \end{figure}
                    
 \textbf{Answer: }318



\textbf{Solution 1}

Note that triangles $\triangle AUV, \triangle BYZ$ and $\triangle CWX$ all have the same height because when they are folded up to create the legs of the table, the top needs to be parallel to the floor. We want to find the maximum possible value of this height, given that no two of $\overline{UV}, \overline{WX}$ and $\overline{YZ}$ intersect inside $\triangle ABC$. Let $h_{A}$ denote the length of the altitude dropped from vertex $A,$ and define $h_{B}$ and $h_{C}$ similarly. Also let $(u, v, w, x, y, z) = (AU, AV, CW, CX, BY, BZ)$. Then by similar triangles
 \begin{align} \frac{u}{AB}=\frac{v}{AC}=\frac{h}{h_{A}}, \\ \frac{w}{CA}=\frac{x}{CB}=\frac{h}{h_{C}}, \\ \frac{y}{BC}=\frac{z}{BA}=\frac{h}{h_{B}}. \end{align}Since $h_{A}=\frac{2K}{23}$ and similarly for $27$ and $30,$ where $K$ is the area of $\triangle ABC,$ we can write
 \begin{align} \frac{u}{30}=\frac{v}{27}=\frac{h}{\frac{2K}{23}}, \\ \frac{w}{27}=\frac{x}{23}=\frac{h}{\frac{2K}{30}}, \\ \frac{y}{23}=\frac{z}{30}=\frac{h}{\frac{2K}{27}}. \end{align}and simplifying gives $u=x=\frac{690h}{2K}, v=y=\frac{621h}{2K}, w=z=\frac{810h}{2K}$. Because no two segments can intersect inside the triangle, we can form the inequalities $v+w\leq 27, x+y\leq 23,$ and $z+u\leq 30$. That is, all three of the inequalities
 \begin{align} \frac{621h+810h}{2K}\leq 27, \\ \frac{690h+621h}{2K}\leq 23, \\ \frac{810h+690h}{2K}\leq 30. \end{align}must hold. Dividing both sides of each equation by the RHS, we have
 \begin{align} \frac{53h}{2K}\leq 1\, \text{since}\, \frac{1431}{27}=53, \\ \frac{57h}{2K}\leq 1\, \text{since}\, \frac{1311}{23}=57, \\ \frac{50h}{2K}\leq 1\, \text{since}\, \frac{1500}{30}=50. \end{align}It is relatively easy to see that $\frac{57h}{2K}\leq 1$ restricts us the most since it cannot hold if the other two do not hold. The largest possible value of $h$ is thus $\frac{2K}{57},$ and note that by Heron's formula the area of $\triangle ABC$ is $20\sqrt{221}$. Then $\frac{2K}{57}=\frac{40\sqrt{221}}{57},$ and the answer is $40+221+57=261+57=\boxed{318}$

 





\textbf{Solution 2}

Note that the area is given by Heron's formula and it is $20\sqrt{221}$. Let $h_i$ denote the length of the altitude dropped from vertex i. It follows that $h_b = \frac{40\sqrt{221}}{27}, h_c  = \frac{40\sqrt{221}}{30}, h_a = \frac{40\sqrt{221}}{23}$. From similar triangles we can see that $\frac{27h}{h_a}+\frac{27h}{h_c} \le 27 \rightarrow h \le \frac{h_ah_c}{h_a+h_c}$. We can see this is true for any combination of a,b,c and thus the minimum of the upper bounds for h yields $h = \frac{40\sqrt{221}}{57} \rightarrow \boxed{318}$.

 





\textbf{Solution 3}

As from above, we can see that the length of the altitude from A is the longest. Thus the highest table is formed when X and Y meet up. Let the distance of this point from B be $x$, making the distance from C $23 - x$. Let $h$ be the height of the table. From similar triangles, we have $\frac{x}{23} = \frac{h}{h_b} = \frac{27h}{2A}$ where A is the area of ABC. Similarly, $\frac{23-x}{23}=\frac{h}{h_c}=\frac{30h}{2A}$. Therefore, $1-\frac{x}{23}=\frac{30h}{2A} \rightarrow1-\frac{27h}{2A}=\frac{30h}{2A}$ and hence $h = \frac{2A}{57} = \frac{40\sqrt{221}}{57}\rightarrow \boxed{318}$.

 




	\item (\textit{AIME I 2011}) Suppose $x$ is in the interval $[0,\pi/2]$ and $\log_{24 \sin x} (24 \cos x) = \frac{3}{2}$.  Find $24 \cot^2 x$.


 
 \textbf{Answer: }192



\textbf{Solution 1}

We can rewrite the given expression as
 \[\sqrt{24^3\sin^3 x}=24\cos x\]Square both sides and divide by $24^2$ to get
 \[24\sin ^3 x=\cos ^2 x\]Rewrite $\cos ^2 x$ as $1-\sin ^2 x$
 \[24\sin ^3 x=1-\sin ^2 x\]
 \[24\sin ^3 x+\sin ^2 x - 1=0\]Testing values using the rational root theorem gives $\sin x=\frac{1}{3}$ as a root, $\sin^{-1} \frac{1}{3}$ does fall in the first quadrant so it satisfies the interval. There are now two ways to finish this problem. 

 \textbf{First way:} Since $\sin x=\frac{1}{3}$, we have
 \[\sin ^2 x=\frac{1}{9}\]Using the Pythagorean Identity gives us $\cos ^2 x=\frac{8}{9}$. Then we use the definition of $\cot ^2 x$ to compute our final answer. $24\cot ^2 x=24\frac{\cos ^2 x}{\sin ^2 x}=24\left(\frac{\frac{8}{9}}{\frac{1}{9}}\right)=24(8)=\boxed{192}$.

 \textbf{Second way:} Multiplying our old equation $24\sin ^3 x=\cos ^2 x$ by $\dfrac{24}{\sin^2x}$ gives 
 \[576\sin x = 24\cot^2x\]So, $24\cot^2x=576\sin x=576\cdot\frac{1}{3}=\boxed{192}$.

 





\textbf{Solution 2}

Like Solution 1, we can rewrite the given expression as
 \[24\sin^3x=\cos^2x\]Divide both sides by $\sin^3x$.
 \[24 = \cot^2x\csc x\]Square both sides.
 \[576 = \cot^4x\csc^2x\]Substitute the identity $\csc^2x = \cot^2x + 1$.
 \[576 = \cot^4x(\cot^2x + 1)\]Let $a = \cot^2x$. Then
 \[576 = a^3 + a^2\].Since $\sqrt[3]{576} \approx 8$, we can easily see that $a = 8$ is a solution. Thus, the answer is $24\cot^2x = 24a = 24 \cdot 8 = \boxed{192}$.

 




	\item (\textit{AIME I 2011}) The probability that a set of three distinct vertices chosen at random from among the vertices of a regular $n$-gon determine an obtuse triangle is $\frac{93}{125}$.  Find the sum of all possible values of $n$.


 
 \textbf{Answer: }503



\textbf{Solution 1}

Inscribe the regular polygon inside a circle. A triangle inside this circle will be obtuse if and only if its three vertices lie on one side of a diameter of the circle. (This is because if an inscribed angle on a circle is obtuse, the arc it spans must be 180 degrees or greater).

 Break up the problem into two cases: an even number of sides $2n$, or an odd number of sides $2n-1$. For polygons with $2n$ sides, the circumdiameter has endpoints on $2$ vertices. There are $n-1$ points on one side of a diameter, plus $1$ of the endpoints of the diameter for a total of $n$ points. For polygons with $2n - 1$ points, the circumdiameter has $1$ endpoint on a vertex and $1$ endpoint on the midpoint of the opposite side. There are also $n - 1$ points on one side of the diameter, plus the vertex for a total of $n$ points on one side of the diameter.

 Case 1: $2n$-sided polygon. There are clearly $\binom{2n}{3}$ different triangles total. To find triangles that meet the criteria, choose the left-most point. There are obviously $2n$ choices for this point. From there, the other two points must be within the $n-1$ points remaining on the same side of the diameter. So our desired probability is$\frac{2n\binom{n-1}{2}}{\binom{2n}{3}}$$=\frac{n(n-1)(n-2)}{\frac{2n(2n-1)(2n-2)}{6}}$$=\frac{6n(n-1)(n-2)}{2n(2n-1)(2n-2)}$$=\frac{3(n-2)}{2(2n-1)}$

 so $\frac{93}{125}=\frac{3(n-2)}{2(2n-1)}$

 $186(2n-1)=375(n-2)$.

 $372n-186=375n-750$

 $3n=564$

 $n=188$ and so the polygon has $376$ sides.

 Case 2: $2n-1$-sided polygon. Similarly, $\binom{2n-1}{3}$ total triangles. Again choose the leftmost point, with $2n-1$ choices. For the other two points, there are again $\binom{n-1}{2}$ possibilities.

 The probability is $\frac{(2n-1)\binom{n-1}{2}}{\binom{2n-1}{3}}$

 $=\frac{3(2n-1)(n-1)(n-2)}{(2n-1)(2n-2)(2n-3)}$

 $=\frac{3(n-2)}{2(2n-3)}$

 so $\frac{93}{125}=\frac{3(n-2)}{2(2n-3)}$

 $186(2n-3)=375(n-2)$

 $375n-750=372n-558$

 $3n=192$

 $n=64$ and our polygon has $127$ sides.

 Adding, $127+376=\boxed{503}$

 





\textbf{Solution 2}

We use casework on the locations of the vertices, if we choose the locations of vertices $v_a, v_b, v_c$ on the n-gon (where the vertices of the n-gon are $v_0, v_1, v_2, ... v_{n-1},$ in clockwise order) to be the vertices of triangle ABC, in order, with the restriction that $a<b<c$.

 By symmetry, we can assume WLOG that the location of vertex A is vertex $v_0$.

 Now, vertex B can be any of $v_1, v_2, ... v_{n-2}$.  We start in on casework.

 Case 1: vertex B is at one of the locations $v_{n-2}, v_{n-3}, ... v_{\lfloor n/2 \rfloor +1}$. (The ceiling function is necessary for the cases in which n is odd.)

 Now, since the clockwise arc from A to B measures more than 180 degrees; for every location of vertex C we can choose in the above restrictions, angle C will be an obtuse angle.

 There are $\lceil n/2 \rceil - 2$ choices for vertex B now (again, the ceiling function is necessary to satisfy both odd and even cases of n).  If vertex B is placed at $v_m$, there are $n - m - 1$ possible places for vertex C.

 Summing over all these possibilities, we obtain that the number of obtuse triangles obtainable from this case is $\frac{(n- \lceil n/2 \rceil - 2)(n - \lceil n/2 \rceil) - 1}{2}$.

 
Case 2: vertex B is at one of the locations not covered in the first case.

 Note that this will result in the same number of obtuse triangles as case 1, but multiplied by 2.  This is because fixing vertex B in $v_0$, then counting up the cases for vertices C, and again for vertices C and A, respectively, is combinatorially equivalent to fixing vertex A at $v_0$, then counting cases for vertex B, as every triangle obtained in this way can be rotated in the n-gon to place vertex A at $v_0$, and will not be congruent to any obtuse triangle obtained in case 1, as there will be a different side opposite the obtuse angle in this case.

 
Therefore, there are $\frac{3(n- \lceil n/2 \rceil - 2)(n - \lceil n/2 \rceil - 1)}{2}$ total obtuse triangles obtainable.

 The total number of triangles obtainable is $1+2+3+...+(n-2) = \frac{(n-2)(n-1)}{2}$.

 The ratio of obtuse triangles obtainable to all triangles obtainable is therefore

 $\frac{\frac{3(n- \lceil n/2 \rceil - 2)(n - \lceil n/2 \rceil - 1)}{2}}{\frac{(n-2)(n-1)}{2}} = \frac{3(n- \lceil n/2 \rceil - 2)(n - \lceil n/2 \rceil - 1)}{(n-2)(n-1)}  = \frac{93}{125}$.

 So, $\frac{(n- \lceil n/2 \rceil - 2)(n - \lceil n/2 \rceil - 1)}{(n-2)(n-1)}  = \frac{31}{125}$.

 Now, we have that $(n-2)(n-1)$ is divisible by $125 = 5^3$.  It is now much easier to perform trial-and-error on possible values of n, because we see that $n \equiv 1,2 \pmod{125}$.

 We find that $n = 127$ and $n = 376$ both work, so the final answer is $127 + 376 = \boxed{503}$.

 





\textbf{Solution 3}

Let the regular $n$-gon be inscribed in a circle, allowing the use of the inscribed angle theorem.We wish to determine the number of three distinct vertices such that they form an obtuse triangle.

 WLOG fix vertex $A$Let $B$ be $k$ edges away from $A$ clockwise	$C$ be $j$ edges away from $A$ in the other direction
                    \begin{figure}[H]
                        \centering
                        \includegraphics[width=0.5\textwidth]{Solution/AIME_I_2011_10_Solution_Figure_1.png}
                    \end{figure}
                    

 To ensure $\triangle ABC$ has no conciding vertices the edges between them $k, j, n - (k+j) > 0$\begin{cases}k > 0\\j > 0\\k + j < n\\\end{cases}The angle of $\triangle ABC$ are $\frac{180^\circ}{n} \cdot (n - k - j), \frac{180^\circ}{n} \cdot j, \frac{180^\circ}{n} \cdot k$ respectively. Solving for the angles $\frac{180^\circ}{n} \cdot x > 90$ for an obtuse,\begin{cases}k > \frac{n}{2}\\j > \frac{n}{2}\\k + j < \frac{n}{2}\end{cases}These are the conditions for each angle of $\triangle ABC$ to be obtuse.Hence we wish to find lattice point $(k, j)$ that satisfy the inequalities.

 As expected these are disjoint, as a triangle contain at most one obtuse angle. 
                    \begin{figure}[H]
                        \centering
                        \includegraphics[width=0.5\textwidth]{Solution/AIME_I_2011_10_Solution_Figure_2.png}
                    \end{figure}
                    We wish to take the union of lattice points within the pink shaded regions. The blue line represents the conditions that make $k$ and $j$ form a non degenerate triangle.

 \textbf{Case} $n = 2c$Focus on of the pink regions. The number of lattice points is $1 + 2 + 3 + \cdots + (c - 2) = \frac{1}{2}(c - 2)(c-1)$Multiply by $3$ for each of the regions resulting in $\frac{3}{2}(c - 2)(c-1)$

 \textbf{Case} $n = 2c - 1$Proceed similarly and get $\frac{3}{2}(c - 2)(c-1)$ too.

 At first the WLOG fixing a single vertex, hence $\times n$Eac $\triangle ABC$ is counted $3$ times, once for each vertex, hence $\div 3$
 \[\frac{3}{2}(c - 2)(c-1) \cdot n \div 3 = \frac{n(c-2)(c-1)}{2}\]Where $c =\left \lfloor \frac{n}{2} \right \rfloor$

 The total number of ways to choose 3 distinct vertices on the $n$-gon is $\binom{n}{3}$

 
 \[\frac{n(c-2)(c-1)}{2 \dbinom{n}{3} } = \frac{93}{125}\]

 \begin{align*}125 \cdot \frac{n(c-2)(c-1)}{2} &= 93 \cdot \dbinom{n}{3}\\ 125 \cdot \frac{n(c-2)(c-1)}{2} &= 93 \cdot \frac{n(n-1)(n-2)}{6}\\125 (c-2)(c-1) &= 31 (n-1)(n-2)\\125 (\left \lfloor \frac{n}{2} \right \rfloor-2)(\left \lfloor \frac{n}{2} \right \rfloor-1) &= 31 (n-1)(n-2)\end{align*}


 \textbf{Case} $n = 2c$\begin{align*}125(c-2)(c-1) &= 31(2c - 1)(2c - 2) \\&= 31 \cdot 2 \cdot (2c - 1)(c-1)\\125(c-2) &= 62(2c-1)\\c &= 188\\n &= 2c = 376\end{align*}


 \textbf{Case} $n = 2c - 1$\begin{align*}125(c-2)(c-1) &= 31(2c - 2)(2c - 3) \\&= 31 \cdot 2 \cdot (c - 1)(2c-3)\\125(c-2) &= 62(2c-3)\\c &= 64\\n &= 2c - 1 = 127\end{align*}


 The final is thus $376 + 127 = \boxed{503}$~clod

 \href{https://mathproblemsolvingskills.wordpress.com/}{MathProblemSolvingSkills.com}

 




	\item (\textit{AIME I 2011}) Let $R$ be the set of all possible remainders when a number of the form $2^n$, $n$ a nonnegative integer, is divided by 1000. Let $S$ be the sum of the elements in $R$. Find the remainder when $S$ is divided by 1000.


 
 \textbf{Answer: }007



\textbf{Solution 1}

Let $k$ be the second smallest integer (since the first one is $3$) such that $2^k \equiv 8 \pmod{125}$. We claim that 
 \[S\equiv 1+2+4+\dots+2^{k-1} \pmod{1000}\] To see this, note that clearly the first three terms can't be congruent to any of the other terms, since they differ mod 8. Then, for $i> j\geq 3$,
 \[2^i\equiv 2^j \pmod{1000} \implies 2^{i-j+3} \equiv 8 \pmod{125}\] contradicting $k$s minimality. Thus, noting that $2^k\equiv 8\pmod{1000}$, 
 \[S\equiv 1+2+4+\dots +2^{k-1} \equiv 2^k-1 \equiv \boxed{007} \pmod{1000}\]

 -sillybone

 





\textbf{Solution 2}

Note that $x \equiv y \pmod{1000} \Leftrightarrow x \equiv y \pmod{125}$ and $x \equiv y \pmod{8}$. So we must find the first two integers 

 




	\item (\textit{AIME I 2011}) Six men and some number of women stand in a line in random order.  Let $p$ be the probability that a group of at least four men stand together in the line, given that every man stands next to at least one other man.  Find the least number of women in the line such that $p$ does not exceed 1 percent.


 
 \textbf{Answer: }594



\textbf{Solution}

Let $n$ be the number of women present, and let \_ be some positive number of women between groups of men.  Since the problem states that every man stands next to another man, there cannot be isolated men.  Thus, there are five cases to consider, where $(k)$ refers to a consecutive group of $k$ men:

 \begin{verbatim}
\_(2)\_(2)\_(2)\_
\end{verbatim}\begin{verbatim}
\_(3)\_(3)\_
\end{verbatim}\begin{verbatim}
\_(2)\_(4)\_
\end{verbatim}\begin{verbatim}
\_(4)\_(2)\_
\end{verbatim}\begin{verbatim}
\_(6)\_
\end{verbatim}For the first case, we can place the three groups of men in between women.  We can think of the groups of men as dividers splitting up the $n$ women.  Since there are $n+1$ possible places to insert the dividers, and we need to choose any three of these locations, we have $\dbinom{n+1}{3}$ ways.

 The second, third, and fourth cases are like the first, only that we need to insert two dividers among the $n+1$ possible locations.  Each gives us $\dbinom{n+1}{2}$ ways, for a total of $3\dbinom{n+1}{2}$ ways.

 The last case gives us $\dbinom{n+1}{1}=n+1$ ways.

 Therefore, the total number of possible ways where there are no isolated men is

 
 \[\dbinom{n+1}{3}+3\dbinom{n+1}{2}+(n+1).\]

 The total number of ways where there is a group of at least four men together is the sum of the third, fourth, and fifth case, or

 
 \[2\dbinom{n+1}{2}+(n+1).\]

 Thus, we want to find the minimum possible value of $n$ where $n$ is a positive integer such that

 
 \[\dfrac{2\dbinom{n+1}{2}+(n+1)}{\dbinom{n+1}{3}+3\dbinom{n+1}{2}+(n+1)}\le\dfrac{1}{100}.\]

 After simplification, we arrive at 
 \[\dfrac{6(n+1)}{n^2+8n+6}\le\dfrac{1}{100}.\]

 Simplifying again, we see that we seek the smallest positive integer value of $n$ such that $n(n-592)\ge594$. Clearly $n>592$, or the left side will not even be positive; we quickly see that $n=593$ is too small but $n=\boxed{594}$ satisfies the inequality.

 




	\item (\textit{AIME I 2011}) A cube with side length 10 is suspended above a plane.  The vertex closest to the plane is labeled $A$. The three vertices adjacent to vertex $A$ are at heights 10, 11, and 12 above the plane.  The distance from vertex $A$ to the plane can be expressed as $\frac{r - \sqrt{s}}{t}$, where $r$, $s$, and $t$ are positive integers. Find $r + s + t$.


 
 \textbf{Answer: }330



\textbf{Solution 1}

Set the cube at the origin with the three vertices along the axes and the plane equal to $ax+by+cz+d=0$, where $a^2+b^2+c^2=1$.  The distance from a point $(X,Y,Z)$ to a plane with equation $Ax+By+Cz+D=0$ is 
 \[\frac{AX+BY+CZ+D}{\sqrt{A^2+B^2+C^2}},\]so the (directed) distance from any point $(x,y,z)$ to the plane is $ax+by+cz+d$.  So, by looking at the three vertices, we have $10a+d=10, 10b+d=11, 10c+d=12$, and by rearranging and summing, 
 \[(10-d)^2+(11-d)^2+(12-d)^2= 100\cdot(a^2+b^2+c^2)=100.\]

 Solving the equation is easier if we substitute $11-d=y$, to get $3y^2+2=100$, or $y=\sqrt {98/3}$.  The distance from the origin to the plane is simply $d$, which is equal to $11-\sqrt{98/3} =(33-\sqrt{294})/3$, so $33+294+3=\boxed{330}$.

 





\textbf{Solution 2}

Let the vertices with distance $10,11,12$ be $B,C,D$, respectively. An equilateral triangle $\triangle BCD$ is formed with side length $10\sqrt{2}$. We care only about the $z$ coordinate: $B=10,C=11,D=12$. It is well known that the centroid of a triangle is the average of the coordinates of its three vertices, so $\text{centroid}=(10+11+12)/3=11$. Designate the midpoint of $BD$ as $M$. Notice that median $CM$ is parallel to the plane because the $\text{centroid}$ and vertex $C$ have the same $z$ coordinate, $11$, and the median contains $C$ and the $\text{centroid}$. We seek the angle $\theta$ of the line:$(1)$ through the centroid $(2)$ perpendicular to the plane formed by $\triangle BCD$, $(3)$ with the plane under the cube. Since the median is parallel to the plane, this orthogonal line is also perpendicular $\textit{in slope}$ to $BD$. Since $BD$ makes a $2-14-10\sqrt{2}$ right triangle, the orthogonal line makes the same right triangle rotated $90^\circ$. Therefore, $\sin\theta=\frac{14}{10\sqrt{2}}=\frac{7\sqrt{2}}{10}$.

 It is also known that the centroid of $\triangle BCD$ is a third of the way between vertex $A$ and $H$, the vertex farthest from the plane. Since $AH$ is a diagonal of the cube, $AH=10\sqrt{3}$. So the distance from the $\text{centroid}$ to $A$ is $10/\sqrt{3}$. So, the $\Delta z$ from $A$ to the centroid is $\frac{10}{\sqrt{3}}\sin\theta=\frac{10}{\sqrt{3}}\left(\frac{7\sqrt{2}}{10}\right)=\frac{7\sqrt{6}}{3}$.

 Thus the distance from $A$ to the plane is $11-\frac{7\sqrt{6}}{3}=\frac{33-7\sqrt{6}}{3}=\frac{33-\sqrt{294}}{3}$, and $33+294+3=\boxed{330}$.

 





\textbf{Solution 3}


                    \begin{figure}[H]
                        \centering
                        \includegraphics[width=0.5\textwidth]{Solution/AIME_I_2011_13_Solution_Figure_1.png}
                    \end{figure}
                    

 





\textbf{Solution 4 (Intuition, not very detailed)}

First, shift the cube down by 11 so the vertices adjacent to $A$ are $-1$, $0$, and $1$ above the plane.

 Consider the two points 1 above and 1 below the plane:
                    \begin{figure}[H]
                        \centering
                        \includegraphics[width=0.5\textwidth]{Solution/AIME_I_2011_13_Solution_Figure_2.png}
                    \end{figure}
                    

 Now we rotate the cube so that all 3 vertices adjacent to $A$ are on the plane. We can calculate $A$ to be $\frac{10\sqrt3}{3}$ below the plane.

 Notice that when $A$ is rotated, it gets slightly closer to the plane.

 
                    \begin{figure}[H]
                        \centering
                        \includegraphics[width=0.5\textwidth]{Solution/AIME_I_2011_13_Solution_Figure_3.png}
                    \end{figure}
                    

 
The triangle shown in the diagram has legs of ration $1:7$, because it is similar to the triangle in the first diagram, since we are rotating. (kind of bad explanation)

 So, we can compute $h$, where $h^2 + \left(\dfrac{h}{7}\right)^2 = \left(\dfrac{10\sqrt3}{3}\right)^2$.

 Solving yields $h=\frac{7\sqrt{6}}{3}$.

 Remember, this is the distance that $A$ is below the plane, so when we shift up by 11, $A$ is $11-\frac{7\sqrt{6}}{3}$ above the plane. Answer extraction yields $\boxed{330}$.

 \href{https://mathproblemsolvingskills.wordpress.com/}{MathProblemSolvingSkills.com}

 


 




	\item (\textit{AIME I 2011}) Let $A_1 A_2 A_3 A_4 A_5 A_6 A_7 A_8$ be a regular octagon.  Let $M_1$, $M_3$, $M_5$, and $M_7$ be the midpoints of sides $\overline{A_1 A_2}$, $\overline{A_3 A_4}$, $\overline{A_5 A_6}$, and $\overline{A_7 A_8}$, respectively.  For $i = 1, 3, 5, 7$, ray $R_i$ is constructed from $M_i$ towards the interior of the octagon such that $R_1 \perp R_3$, $R_3 \perp R_5$, $R_5 \perp R_7$, and $R_7 \perp R_1$.  Pairs of rays $R_1$ and $R_3$, $R_3$ and $R_5$, $R_5$ and $R_7$, and $R_7$ and $R_1$ meet at $B_1$, $B_3$, $B_5$, $B_7$ respectively.  If $B_1 B_3 = A_1 A_2$, then $\cos 2 \angle A_3 M_3 B_1$ can be written in the form $m - \sqrt{n}$, where $m$ and $n$ are positive integers.  Find $m + n$.


 
 \textbf{Answer: }037



\textbf{Solution 1}

We use coordinates. Let the octagon have side length $2$ and center $(0, 0)$. Then all of its vertices have the form $(\pm 1, \pm\left(1+\sqrt{2}\right))$ or $(\pm\left(1+\sqrt{2}\right), \pm 1)$.

 By symmetry, $B_{1}B_{3}B_{5}B_{7}$ is a square. Thus lines $\overleftrightarrow{B_{1}B_{3}}$ and $\overleftrightarrow{B_{5}B_{7}}$ are parallel, and its side length is the distance between these two lines. However, this is given to be the side length of the octagon, or $2$.

 Suppose the common slope of the lines is $m$ and let $m=\tan\theta$. Then, we want to find 
 \[\cos 2\left(90-\theta\right)=2\cos^{2}\left(90-\theta\right)-1=2\sin^{2}\theta-1.\]

 It can easily be seen that the equations of the lines are 
 \begin{align*} B_{1}B_{3}: y-mx+m\left(1+\sqrt{2}\right)=0 \\ B_{5}B_{7}: y-mx-m\left(1+\sqrt{2}\right)=0.\end{align*}
 By the \href{https://www.cuemath.com/geometry/distance-between-two-lines/}{distance between parallel lines formula}, a corollary of the \href{https://brilliant.org/wiki/dot-product-distance-between-point-and-a-line/}{point to line distance formula}, the distance between these two lines is 
 \[\frac{|c_{2}-c_{1}|}{\sqrt{a^{2}+b^{2}}}=\frac{2m\left(1+\sqrt{2}\right)}{\sqrt{m^{2}+1}}.\] Since we want this to equal $2$, we have 
 \begin{align*}\frac{2m\left(1+\sqrt{2}\right)}{\sqrt{m^{2}+1}}&=2 \\ 4m^{2}\left(3+2\sqrt{2}\right)&=4m^{2}+4 \\ \left(12+8\sqrt{2}\right)m^{2}&=4m^{2}+4 \\ \left(8+8\sqrt{2}\right)m^{2}&=4 \\ m^{2}&=\frac{4}{8+8\sqrt{2}} \\ \Rightarrow m^{2}=\tan^{2}\theta=\frac{\sin^{2}\theta}{\cos^{2}\theta}&=\frac{1}{2+2\sqrt{2}}.\end{align*}
 Since $\sin^{2}\theta+\cos^{2}\theta=1,$ we have $\sin^{2}\theta=\frac{1}{3+2\sqrt{2}}$. Thus 
 \[2\sin^{2}\theta-1=\frac{2}{3+2\sqrt{2}}-1=\frac{-1-2\sqrt{2}}{3+2\sqrt{2}}=\frac{\left(-1-2\sqrt{2}\right)\left(3-2\sqrt{2}\right)}{\left(3+2\sqrt{2}\right)\left(3-2\sqrt{2}\right)}=\frac{5-4\sqrt{2}}{1}=5-\sqrt{32}.\] The answer is $\boxed{037}$.

 


 





\textbf{Solution 2}

Let $\theta=\angle M_1 M_3 B_1$. Thus we have that $\cos 2 \angle A_3 M_3 B_1=\cos \left(2\theta + \frac{\pi}{2} \right)=-\sin2\theta$.

 Since $A_1 A_2 A_3 A_4 A_5 A_6 A_7 A_8$ is a regular octagon and $B_1 B_3 = A_1 A_2$, let $k=A_1 A_2 = A_2 A_3 = B_1 B_3$.

 
Extend $\overline{A_1 A_2}$ and $\overline{A_3 A_4}$ until they intersect. Denote their intersection as $I_1$. Through similar triangles & the $45-45-90$ triangles formed, we find that $M_1 M_3=\frac{k}{2}(2+\sqrt2)$.

 We also have that $\triangle M_7 B_7 M_1 =\triangle M_1 B_1 M_3$ through ASA congruence ($\angle B_7 M_7 M_1 =\angle B_1 M_1 M_3$, $M_7 M_1 = M_1 M_3$, $\angle B_7 M_1 M_7 =\angle B_1 M_3 M_1$). Therefore, we may let $n=M_1 B_7 = M_3 B_1$.

 Thus, we have that $\sin\theta=\frac{n-k}{\frac{k}{2}(2+\sqrt2)}$ and that $\cos\theta=\frac{n}{\frac{k}{2}(2+\sqrt2)}$. Therefore $\cos\theta-\sin\theta=\frac{k}{\frac{k}{2}(2+\sqrt2)}=\frac{2}{2+\sqrt2}=2-\sqrt2$.

 Squaring gives that $\sin^2\theta - 2\sin\theta\cos\theta + \cos^2\theta = 6-4\sqrt2$ and consequently that $-2\sin\theta\cos\theta = 5-4\sqrt2 = -\sin2\theta$ through the identities $\sin^2\theta + \cos^2\theta = 1$ and $\sin2\theta = 2\sin\theta\cos\theta$.

 Thus we have that $\cos 2 \angle A_3 M_3 B_1=5-4\sqrt2=5-\sqrt{32}$. Therefore $m+n=5+32=\boxed{037}$.

 





\textbf{Solution 3}

Let $A_1A_2 = 2$.  Then $B_1$ and $B_3$ are the projections of $M_1$ and $M_5$ onto the line $B_1B_3$, so $2=B_1B_3=-M_1M_5\cos x$, where $x = \angle A_3M_3B_1$.  Then since $M_1M_5 = 2+2\sqrt{2}, \cos x = \dfrac{-2}{2+2\sqrt{2}}= 1-\sqrt{2}$, $\cos 2x = 2\cos^2 x -1 = 5 - 4\sqrt{2} = 5-\sqrt{32}$, and $m+n=\boxed{037}$.

 





\textbf{Solution 4}

Notice that $R_3$ and $R_7$ are parallel ($B_1B_3B_5B_7$ is a square by symmetry and since the rays are perpendicular) and $B_1B_3=B_3B_5=s=$ the distance between the parallel rays. If the regular hexagon as a side length of $s$, then $M_3M_7$ has a length of $s+s\sqrt{2}$. Let $X$ be on $R_3$ such that $M_7X$ is perpendicular to $M_3X$, and $\phi=\angle M_7M_3X$. The distance between $R_3$ and $R_7$ is $s=M_7X$, so $\sin\phi=\frac{s}{s+s\sqrt{2}}=\frac{1}{1+\sqrt{2}}$. 

 Since we are considering a regular hexagon, $M_3$ is directly opposite to $M_7$ and $\angle A_3M_3B_1=90 ^\circ +\phi$. All that's left is to calculate $\cos 2\angle A_3M_3B_1=\cos^2(90^\circ+\phi)-\sin^2(90^\circ+\phi)=\sin^2\phi-\cos^2\phi$. By drawing a right triangle or using the Pythagorean identity, $\cos^2\phi=\frac{2+2\sqrt2}{3+2\sqrt2}$ and $\cos 2\angle A_3M_3B_1=\frac{-1-2\sqrt2}{3+2\sqrt2}=5-4\sqrt2=5-\sqrt{32}$, so $m+n=\boxed{037}$.

 


 





\textbf{Solution 5}

Assume that $A_1A_2=1.$Denote the center $O$, and the midpoint of $B_1$ and $B_3$ as $B_2$. Then we have that
 \[\cos\angle A_3M_3B_1=\cos(\angle A_3M_3O+\angle OM_3B_1)=-\sin(\angle OM_3B_1)=-\frac{OB_2}{OM_3}=-\frac{1/2}{1/2+\sqrt2/2}=-\frac{1}{\sqrt2+1}=1-\sqrt2.\]Thus, by the cosine double-angle theorem,
 \[\cos2\angle A_3M_3B_1=2(1-\sqrt2)^2-1=5-\sqrt{32},\]so $m+n=\boxed{037}$.

 


 \href{https://mathproblemsolvingskills.wordpress.com/}{MathProblemSolvingSkills.com}

 


 


 




	\item (\textit{AIME I 2011}) For some integer $m$, the polynomial $x^3 - 2011x + m$ has the three integer roots $a$, $b$, and $c$.  Find $|a| + |b| + |c|$.


 
 \textbf{Answer: }098



\textbf{Solution 1}

From Vieta's formulas, we know that $a+b+c = 0$, and $ab+bc+ac = -2011$. Thus $a = -(b+c)$. All three of $a$, $b$, and $c$ are non-zero: say, if $a=0$, then $b=-c=\pm\sqrt{2011}$ (which is not an integer). $\textsc{wlog}$, let $|a| \ge |b| \ge |c|$. If $a > 0$, then $b,c < 0$ and if $a < 0$, then $b,c > 0,$ from the fact that $a+b+c=0$. We have 
 \[-2011=ab+bc+ac = a(b+c)+bc = -a^2+bc\]Thus $a^2 = 2011 + bc$. We know that $b$, $c$ have the same sign, so product $bc$ is always positive. So $|a| \ge 45 = \lceil \sqrt{2011} \rceil$. 

 Also, if we fix $a$, $b+c$ is fixed, so $bc$ is maximized when $b = c$ . Hence, 
 \[2011 = a^2 - bc > \frac{3}{4}a^2 \qquad \Longrightarrow \qquad a ^2 < \frac{4}{3}\cdot 2011 = 2681+\frac{1}{3}\]So $|a| \le 51$. Thus we have bounded $a$ as $45\le |a| \le 51$, i.e. $45\le |b+c| \le 51$ since $a=-(b+c)$. Let's analyze $bc=(b+c)^2-2011$. Here is a table: 

 <table border="1" cellpadding="5" cellspacing="0" style="text-align:center;"><tbody><tr><th>$|a|$</th><th>$bc=a^2-2011$</th></tr><tr><td>$45$</td><td>$14$</td></tr><tr><td>$46$</td><td>$105$</td></tr><tr><td>$47$</td><td>$198$</td></tr><tr><td>$48$</td><td>$293$</td></tr><tr><td>$49$</td><td>$390$</td></tr></tbody></table>


 We can tell we don't need to bother with $45$, 

 $105 = (3)(5)(7)$, So $46$ won't work. $198/47 > 4$, 

 $198$ is not divisible by $5$, $198/6 = 33$, which is too small to get $47$.

 $293/48 > 6$, $293$ is not divisible by $7$ or $8$ or $9$, we can clearly tell that $10$ is too much.

 
Hence, $|a| = 49$, $a^2 -2011 = 390$. $b = 39$, $c = 10$.

 Answer: $\boxed{098}$

 





\textbf{Solution 2}

Starting off like the previous solution, we know that $a + b + c = 0$, and $ab + bc + ac = -2011$.

 Therefore, $c = -b-a$.

 Substituting, $ab + b(-b-a) + a(-b-a) = ab-b^2-ab-ab-a^2 = -2011$.

 Factoring the perfect square, we get: $ab-(b+a)^2=-2011$ or $(b+a)^2-ab=2011$.

 Therefore, a sum ($a+b$) squared minus a product ($ab$) gives $2011$..

 


 We can guess and check different $a+b$’s starting with $45$ since $44^2 < 2011$.

 $45^2 = 2025$ therefore $ab = 2025-2011 = 14$. 

 Since no factors of $14$ can sum to $45$ ($1+14$ being the largest sum), a + b cannot equal $45$.

 $46^2 = 2116$ making $ab = 105 = 3 * 5 * 7$.

 $5 * 7 + 3 < 46$ and $3 * 5 * 7 > 46$ so $46$ cannot work either.

 


 We can continue to do this until we reach $49$.

 $49^2 =  2401$ making $ab = 390 = 2 * 3 * 5* 13$.

 $3 * 13 + 2* 5 = 49$, so one root is $10$ and another is $39$. The roots sum to zero, so the last root must be $-49$.

 


 $|-49|+10+39 = \boxed{098}$.

 





\textbf{Solution 3}

Let us first note the obvious that is derived from Vieta's formulas: $a+b+c=0, ab+bc+ac=-2011$. Now, due to the first equation, let us say that $a+b=-c$, meaning that $a,b>0$ and $c<0$. Now, since both $a$ and $b$ are greater than 0, their absolute values are both equal to $a$ and $b$, respectively. Since $c$ is less than 0, it equals $-a-b$. Therefore, $|c|=|-a-b|=a+b$, meaning $|a|+|b|+|c|=2(a+b)$. We now apply Newton's sums to get that $a^2+b^2+ab=2011$,or $(a+b)^2-ab=2011$. Solving, we find that $49^2-390$ satisfies this, meaning $a+b=49$, so $2(a+b)=\boxed{098}$.

 





\textbf{Solution 4 (Quadratic)}

From \href{https://artofproblemsolving.com/wiki/index.php?title=Vieta\%27s\_Formulas}{Vieta's Formulas} we have $a + b + c = 0$ and $a(b + c) + bc = -2011.$ $b + c = -a \implies -a^2 + bc = -2011.$ $a^2 = (b + c)^2 \implies bc + 2011 = (b + c)^2.$  So now we have a simple looking two-variable quadratic equation. From here, we can solve for $c$ in terms of $b$ using the \href{https://artofproblemsolving.com/wiki/index.php?title=Quadratic\_formula}{quadratic formula} and see if we can do something with the discriminant. $bc + 2011 = b^2 + 2bc + c^2.$ $c^2 + (b)c + (b^2 - 2011) = 0.$ So $c = \frac{-b \pm \sqrt{8044 - 3b^2}}{2}.$ So $8044 - 3b^2$ must be a perfect square. $8044$ is around $89^2$, so we can start from here and work downwards. Immediately, we see that if $8044 - 3b^2 = 88^2, b = 10.$ Will this work? If $b = 10,$ then $c$ can be $-49$ or $39$. So we have two cases: 

 
\textbf{Case 1:} $c = -49$ 

 $b + c = -39$ and $bc = -490$ 

 $a(-39) - 490 = -2011 \implies a = -39.$ But this doesn't satisfy the equation $a + b + c = 0,$ so this case won't work.

 
\textbf{Case 2:} $c = 39$ 

 $b + c = 49$ and $bc = 390$ 

 $a(49) + 390 = -2011 \implies a = -49.$ Since $a + b + c = 10 + 39 - 49 = 0$ this works.

 So our answer is $10 + 39 + 49 = \boxed{98}.$

 





\textbf{Solution 5 (mod to help bash)}

First, derive the equations $a=-b-c$ and $ab+bc+ca=-2011\implies b^2+bc+c^2=2011$. Since the product is negative, $a$ is negative, and $b$ and $c$ positive. Now, a simple mod 3 testing of all cases shows that $b\equiv \{1,2\} \pmod{3}$, and $c$ has the repective value. We can choose $b$ not congruent to 0, make sure you see why. Now, we bash on values of $b$, testing the quadratic function to see if $c$ is positive. You can also use a delta argument like solution 4, but this is simpler. We get that for $b=10$, $c=39, -49$. Choosing $c$ positive we get $a=-49$, so $|a|+|b|+|c|=10+49+39=\boxed{098}$~firebolt360

 





\textbf{Solution 6}

Note that $-c=b+a$, so $c^2=a^2+2ab+b^2$, or $-c^2+ab=-a^2-ab-b^2$. Also, $ab+bc+ca=-2011$, so $(a+b)c+ab=-c^2+ab=-2011$. Substituting $-c^2+ab=-a^2-ab-b^2$, we can obtain $a^2+ab+b^2=2011$, or $\frac{a^3-b^3}{a-b}=2011$. If it is not known that $2011$ is prime, it may be proved in $5$ minutes or so by checking all primes up to $43$. If $2011$ divided either of $a, b$, then in order for $a^3-b^3$ to contain an extra copy of $2011$, both $a, b$ would need to be divisible by $2021$. But then $c$ would also be divisible by $2011$, and the sum $ab+bc+ca$ would clearly be divisible by $2011^2$. 

 By LTE, $v_{2011}(a^3-b^3)=v_{2011}(a-b)$ if $a-b$ is divisible by $2011$ and neither $a,b$ are divisible by $2011$. Thus, the only possibility remaining is if $a-b$ did not divide $2011$. Let $a=k+b$. Then, we have $(b+k)^3-b^3=2011k$. Rearranging gives $3b(b+k)=2011-k^2$. As in the above solutions, we may eliminate certain values of $k$ by using mods. Then, we may test values until we obtain $k=29$, and $a=10$. Thus, $b=39$, $c=-49$, and our answer is $49+39+10=098$.

 





\textbf{Solution 7}

For the monic cubic $x^3 - 2011x + m$ with integer roots $a,b,c$, we have

 $a+b+c=0,\quad ab+ac+bc=-2011,\quad abc=-m.$

 Set $c=-(a+b).$ Then 

 
 \[ab + a(-a-b) + b(-a-b) = -2011,\]

 which simplifies to 

 
 \[-a^2 - ab - b^2 = -2011 \;\;\Longrightarrow\;\; a^2+ab+b^2=2011.\]

 Treat this as a quadratic in $a$. Its discriminant must be a perfect square:

 
 \[\Delta = b^2 - 4(b^2 - 2011) = 8044 - 3b^2 = k^2 \quad \text{for some integer }k.\]

 Thus $3b^2 + k^2 = 8044.$ Checking integer $b$ with $|b|\leq 51$ (since $3b^2 \leq 8044$) yields solutions only when  $b=10,\;39,\;49 \quad (\text{up to sign}),$ which produce the root triples (up to permutation)

 
 \[(39,\,10,\,-49)\]

 and the sign-permuted equivalent

 $(-39,\,-10,\,49).$

 These satisfy $a+b+c=0 \quad \text{and} \quad ab+ac+bc=-2011.$ Therefore $|a|+|b|+|c| = 39+10+49 =\boxed{98}$.

 




\end{enumerate}